This section provides a practical guide for converting between PP and S in
court filings, statutory compliance documents, and expert reports.

\subsection{When to Report S vs.\ PP}

\begin{itemize}
  \item \textbf{Report PP} when: filing in federal court; using the Princeton
    Gerrymandering Project's grading framework; comparing to prior expert reports
    that use PP; or when the audience is familiar with $[0,1]$ scales.
  \item \textbf{Report S} when: filing with Colorado's IRC; statutory context
    requires $S \leq c$ language; audience is a legislative body more comfortable
    with ``distance from circular'' framing.
  \item \textbf{Report both} when: filing in a jurisdiction that has used both
    conventions, or when the report is intended for a general audience that
    includes both court and legislative readers.
\end{itemize}

\subsection{Example Expert Report Language}

\begin{quote}
  \emph{The proposed plan achieves a mean Polsby-Popper score of $0.22$ across
  all 14 NC congressional districts, equivalent to a mean Schwartzberg score of
  $S = 1/\sqrt{0.22} \approx 2.13$. Individual district PP values range from
  $0.13$ to $0.36$ (S $\in [1.67, 2.77]$). These values substantially exceed
  the NC enacted 2020 map's reported mean PP of approximately $0.12$--$0.15$
  (S $\approx 2.58$--$2.89$).}
\end{quote}

\subsection{Conversion for Statutory Compliance}

When a statute mandates $S \leq c$, use the equivalence $S \leq c \Leftrightarrow \text{PP} \geq 1/c^2$.
For Colorado's informal $S \leq 2.0$ benchmark:
\[
  \text{PP} \geq \frac{1}{(2.0)^2} = 0.25.
\]
The bisect command:
\begin{verbatim}
bisect label-analyze <label> --year 2020 --types schwartzberg
\end{verbatim}
reports both S values and the equivalent PP for each district,
with a statutory compliance flag if any district exceeds the specified threshold.
